%======================================================================
% CardaVI: project note
%======================================================================
\documentclass[a4paper, 11pt]{article}
\usepackage[utf8]{inputenc}
\usepackage[T1]{fontenc}
\usepackage{amsmath, amssymb, amsthm}
\usepackage[margin=1in]{geometry}
\usepackage{enumitem}
\usepackage{booktabs}
\usepackage{hyperref}
\hypersetup{colorlinks=true, linkcolor=blue, citecolor=blue, urlcolor=blue}

\DeclareMathOperator*{\argmin}{arg\,min}
\newcommand{\E}{\mathbb{E}}
\newcommand{\CE}{\mathrm{CE}}
\newcommand{\NB}{\mathrm{NB}}
\newcommand{\nnz}{\mathrm{nnz}}
\newcommand{\klow}{k^{\mathrm{low}}}
\newcommand{\khigh}{k^{\mathrm{high}}}
\newcommand{\tofill}[1]{\par\smallskip\noindent\emph{[To be completed: #1]}\par\smallskip}

\theoremstyle{plain}
\newtheorem{prop}{Proposition}
\theoremstyle{remark}
\newtheorem{rem}{Remark}

\title{CardaVI\\[0.3em] \large Mechanistic gene regulatory network inference at the scale of the transcriptome:\\ a regulatory core, hidden regulators and a single objective}
\author{Project note}
\date{\today}

\begin{document}
\maketitle

%======================================================================
\section*{Project note}
%======================================================================

\paragraph{Objective.} CardamomOT~\cite{mauge2026cardamomot} infers a gene regulatory network (GRN), hidden protein trajectories and kinetic rates from scRNA-seq time series, and gives a generative model of the cells. Without strong constraints on the possible interactions, it is limited to small gene sets selected before the inference. The question of the project is whether causal regulation can be inferred at the scale of the transcriptome from snapshots, including for regulators that are never measured. The latent state of a cell is the set of proteins of a small core of regulators, completed by a few hidden variables: it follows the mechanistic dynamics of CardamomOT and predicts the law of every gene. The project has two stages. The first extends CardamomOT directly: the couplings use all the genes, and the core is selected within the inference among candidate regulators given by a prior. The second, CardaVI, makes every step of the inference descend a single objective, with an encoder that reads the hidden variables in each cell.

\paragraph{The thread.} The number of genes can grow; the number of regulators that the data can tell apart cannot (Section~\ref{sec:cap}). Most genes should therefore be targets of the core, not free regulators. They then become an information source: through the coupling, thousands of targets constrain the proteins of the core. The coupling is also what keeps the problem constrained: once an ancestor is chosen, it gives one equation per gene and per cell, where a fit on marginals alone is under-constrained (Section~\ref{sec:objective}). The two stages share this objective, written from the start, and differ in how it is optimised.

\paragraph{What already exists in the code.} The negative binomial mixture, the rates, the mechanistic bridges, the OT cost with its network term, the recomputation of the end states on the pairs drawn from the coupling only, the fits per target with the QUT penalty, the update of the basins, and the free-running fit of the simulation model with reparametrised Gamma bursts (\texttt{inference/sde.py}), which gives the dynamics of the hidden variables. The current gene selection gives an initial core.

\paragraph{Where we stand.} Nothing specific to this project is implemented. Propositions~\ref{prop:path} and~\ref{prop:bound} and the transition law~\eqref{eq:transition} have been checked numerically on small examples. The claims of Section~\ref{sec:cap} on the effective rank have not been measured on our datasets, and the central hypothesis, that genes outside the core improve the network of the core, has not been tested.

\paragraph{Next steps.} In the order of Section~\ref{sec:validation}, each step deciding whether the next one is worth doing:
\begin{enumerate}[nosep]
\item Simulated data with many genes and a known core: all the genes and a prior with decoys, against CardamomOT on the right genes only.
\item Lock-in: remove a true regulator from the initial core, and check whether it comes back.
\item Does the recovery of the core improve with the number of genes outside the core?
\item The objective along the iterations of the first stage: which steps do not descend it?
\item A hidden regulator: remove the mRNA of a known regulator, and recover its protein and its regulators.
\end{enumerate}

%======================================================================
\section{Introduction}\label{sec:intro}
%======================================================================

Generative models of single-cell data scale to the whole transcriptome by working in a low-dimensional latent space. In scVI~\cite{lopez2018deep}, the latent variable has a Gaussian prior and the decoder is a free neural network. biVI~\cite{carilli2024biophysical} replaces the negative binomial of scVI by the distributions of biophysical models of transcription, so that the decoder outputs interpretable kinetic parameters. In both, the cells are independent draws at stationarity: there is no time, no network, and no feedback of the gene products on the latent state.

CardamomOT has all three: the proteins are produced by the genes and regulate them through a GRN, along trajectories reconstructed by mechanistic optimal transport. It works in the space of the genes, with one protein per gene. The present project goes one step beyond biVI in the direction of CardamomOT: the latent variable is the state of the regulators and follows their mechanistic dynamics, the decoder is the GRN, and the law of each gene depends on the recent path of its burst frequency. Models with a network at scale exist (FLeCS~\cite{bertin2025scalable} uses coupled differential equations structured by a gene network), as do models that infer an unobserved regulator from its targets (latent force models~\cite{moss2020gene}, on small bulk time series). To our knowledge, none combines bursty stochastic dynamics, hidden regulators and transcriptome-scale targets on snapshot data.

%======================================================================
\section{Background: CardamomOT}\label{sec:background}
%======================================================================

Each gene $g$ is transcribed in bursts at a frequency set by the proteins through the GRN~\cite{herbachinferring2017}. Counts are modelled, given a level $k$ of the burst frequency, by $x_g \sim \NB(k, c_g/s)$, of shape $k$ and rate $c_g/s$, with $s$ the depth of the cell. The burst size is $b_g = 1/c_g$, and the levels are bounded by $\klow_g$ and $\khigh_g$, estimated gene by gene with the mixture. The degradation rates $d_{0,g}$ and $d_{1,g}$ come from tables. Regulation enters through the sigmoid
\[
k_g(P) = \klow_g + (\khigh_g - \klow_g)\,\sigma_g(P), \qquad \sigma_g(P) = \mathrm{logistic}\Big(\beta_g + \sum_r \theta_{rg} P_r\Big),
\]
written here in its logistic form; the other regulation functions of the model learning apply in the same way.

The inference alternates three steps. Trajectories: given the network, each trajectory state $i$ at time $t_a$ is coupled to the real cells $j$ at time $t_b$, with $\Delta = t_b - t_a$, by an entropic OT of cost
\begin{equation}\label{eq:current-cost}
c_{ij} = w A_{ij} + (1 - w)\Big[\big(1 - \tfrac{1}{G}\big) B_{ij} + \tfrac{1}{G} D_{ij}\Big],
\qquad B_{ij} = \sum_g w_g \CE\big(\sigma_g(P_{ij}), \tau_{gj}\big),
\end{equation}
where $P_{ij}$ is the end state of the mechanistic bridge from state $i$ with the modes of $i$ and $j$, $\tau_{gj} \in [0,1]$ is the target of gene $g$ in cell $j$ given by its modes, $A_{ij}$ is a distance on the mRNA whose weight $w$ decreases along the iterations, and $D_{ij}$ a distance on the proteins. One target is then drawn for each state, and the end states are recomputed for the drawn pairs only. Network: the regression of the targets on the proteins, target by target. Basins: the update of the mode of each cell and gene given the network. A free-running fit of the simulation model then refines the network and the rates.

The computing time grows as $n_1 n_2 G^2$ per batch of $n_1$ states and $n_2$ cells, since each pair requires the bridges of $G$ genes and the sigmoid of $G$ targets on $G$ proteins.

%======================================================================
\section{Why the number of regulators is capped}\label{sec:cap}
%======================================================================

The regression of a target $g$ on the proteins has a design matrix $X$ of cells by regulators. Cells concentrate around a few basins and along low-dimensional transients, so $X$ has an effective rank $K$ of the order of the intrinsic dimension of the trajectories. Beyond it, $\theta_g + v$ with $Xv \approx 0$ fits as well as $\theta_g$: the sparsity penalty picks one solution among collinear regulators, and this choice is arbitrary. Only the projection of $\theta_g$ on $K$ directions is identifiable, whatever the number of genes. The transients and the free-running fit add information, within the same dimension. From marginals alone, part of the dynamics is not identifiable anyway~\cite{weinreb2018fundamental, lavenant2021towards}.

\tofill{effective rank of $X$ on our datasets, and stability of the edges between collinear regulators.}

%======================================================================
\section{The framework}\label{sec:framework}
%======================================================================

\subsection{The latent state}

The latent state of a cell is $z = (P, z^{\mathrm{h}})$: the proteins $P$ of the $K$ regulators of the core, and the states $z^{\mathrm{h}}$ of a few hidden variables. The stimuli are exogenous components of $z$, set by their schedules as in CardamomOT.

Regulators are of two kinds. An \emph{anchored} regulator is a gene: its mRNA is observed, its levels come from its mixture and its rates from the tables. Its protein is latent, as in the current model. Not every regulator needs to be a transcription factor: a chromatin regulator, for instance, can be anchored. A \emph{hidden} variable has no observed mRNA (Section~\ref{sec:dynamics}). It is part of $z$ and cannot be deduced from the core: the core only predicts its production, and its state is read on its targets. A hidden variable whose state were a function of the core would only add a delay.

The latent state predicts a law, not an expression. The recent path of $z$ gives the burst frequency of every gene (Section~\ref{sec:decoder}), and the counts remain stochastic given it. All the correlation between genes must therefore go through $z$, which gives a testable criterion: residual correlations between genes given $z$ show that something is missing from $z$.

The framework has five components, which also receive the extensions of CardamomOT (Section~\ref{sec:extensions}): the dynamics of $z$, the decoder from $z$ to the genes, the measurements of $z$, the constraints on the coupling, and the prior on $\theta$.

\subsection{Decoder}\label{sec:decoder}

Only the $K$ rows of $\theta$ (regulator to target) of the core are non-zero. The other genes are \emph{satellites}: targets without outgoing edges, which need neither proteins nor bridges. Their parameters are their rows of $\theta$ restricted to the core, fitted target by target with only $K$ regressors.

At fixed burst frequency, the NB is the stationary law of the bursty model, as in biVI. Given the \emph{path} of the frequency, the law of the mRNA is explicit.

\begin{prop}\label{prop:path}
Let the mRNA of gene $g$ receive exponential bursts of mean $b$ at rate $d_0\, k(t)$ and decay at rate $d_0$. Then $M(t)$ is a shot noise with
\[
\log \E\big[e^{-\lambda M(t)}\big] = - d_0 \int_0^\infty k(t-u)\, \frac{\lambda b\, e^{-d_0 u}}{1 + \lambda b\, e^{-d_0 u}}\, du,
\]
and its mean and variance are $b\,\bar k(t)$ and $b^2\, \tilde k(t)$, where
\[
\bar k(t) = \int_0^\infty d_0\, e^{-d_0 u}\, k(t-u)\, du, \qquad
\tilde k(t) = \int_0^\infty 2 d_0\, e^{-2 d_0 u}\, k(t-u)\, du .
\]
At constant $k$, $M(t)$ follows the Gamma law of shape $k$ and scale $b$.
\end{prop}

\begin{proof}
Campbell's formula for the Poisson process of the bursts, with $\E[e^{-\lambda X}] = (1 + \lambda a)^{-1}$ for an exponential mark of mean $a = b\, e^{-d_0 u}$. The substitution $v = \lambda b\, e^{-d_0 u}$ gives $(1 + \lambda b)^{-k}$ at constant $k$.
\end{proof}

Matching the first two moments with a Gamma law gives the law of the counts along a path of $z$:
\begin{equation}\label{eq:decoder}
x_{gj} \mid z \sim \NB\Big(\text{shape } \frac{\bar k_g^2}{\tilde k_g},\ \text{mean } s_j\, b_g\, \bar k_g\Big), \qquad
k_g(\tau) = \klow_g + (\khigh_g - \klow_g)\,\sigma_g\big(z(\tau)\big),
\end{equation}
whose integrals are quadratures along the recent path of $z$, like the nodes of the mRNA memory of CardamomOT. It is an approximation: on a recent switch-on, its lower tail is visibly off (first decile $0.97$ against $0.72$ in a Monte Carlo simulation of the shot noise, with the parameters of the remark below), while the mean, variance and upper quantiles match. This law is a result on the model, for its simulations; the coupling keeps the cross-entropy on the modes (Section~\ref{sec:objective}).

\begin{rem}[A snapshot sees the direction of a switch]
After a switch-on, $\tilde k > \bar k$ since $\tilde k$ weights the recent past more; after a switch-off, $\tilde k < \bar k$. Two populations with the same mean thus differ in dispersion according to the direction of the switch. With $d_0 = 1$, $b = 2$, levels $0.5$ and $4$ and a switch $0.7$ time units ago, the means are $4.52$ (on) and $4.48$ (off), the variances $12.5$ and $5.5$. More generally, each mRNA filters the frequency with its own $d_0$, so that one cell contains a window of its recent past, provided that $d_0$ is fixed in advance. Unspliced and spliced counts, the data of biVI, would add to it naturally. Whether this signal survives in real data is an open question.
\end{rem}

\subsection{Dynamics}\label{sec:dynamics}

The proteins of the anchored regulators follow the dynamics of CardamomOT: bridges between the modes of their mRNA, with dilution and growth when they are on. A hidden variable has no state for its mRNA, since its $d_0$ and its burst size could not be identified. It follows directly the protein-only model of the free-running fit (mode \texttt{sde1}):
\begin{equation}\label{eq:hidden}
dz_h = d_{1,h}\big(\sigma_h(z) - z_h\big)\, dt + \sqrt{2\, \sigma_h(z)\, d_{1,h}\, a_h}\, dW,
\end{equation}
written in its diffusion limit (the simulations use Gamma bursts). The sigmoid $\sigma_h$ plays the role of the mRNA: it is the level towards which the protein relaxes. The parameters are the incoming edges $\Theta_h$, the rate $d_{1,h}$ and a single noise parameter $a_h$ (the ratio $\varepsilon / k_1$ of the code). The levels are a convention of scale, since an affine transformation of $z_h$ is absorbed by the $\beta$ and the edges of its targets: we set $z_h \in [0,1]$.

At frozen $\sigma_h$, the transition over an interval $\Delta$ is Gaussian in~\eqref{eq:hidden}, with
\begin{equation}\label{eq:transition}
m_h = e^{-d_{1,h}\Delta}\, z_h(i) + \big(1 - e^{-d_{1,h}\Delta}\big)\, \sigma_h, \qquad
v_h = a_h\, \sigma_h\, \big(1 - e^{-2 d_{1,h}\Delta}\big).
\end{equation}
An mRNA of the hidden variable only exists implicitly, per pair of cells, given an ancestor: the mean production that brings the state from the ancestor to the cell,
\begin{equation}\label{eq:implied}
\widehat M_h(i,j) = \frac{z_h(j) - e^{-d_{1,h}\Delta}\, z_h(i)}{1 - e^{-d_{1,h}\Delta}} .
\end{equation}
Comparing it with $\sigma_h$ is the reachability term of Section~\ref{sec:objective}, up to a factor, without the division that amplifies noise when $d_{1,h}\Delta$ is small.

Hidden variables can represent two different things: an unobserved regulator of the network, with a small $a_h$, and a nuisance factor such as the cell cycle, with a large $a_h$, which behaves as a static factor per cell. The same parameter separates them, without two kinds of latent variables. The rate $d_{1,h}$ is only identifiable if $d_{1,h}\Delta$ is of the order of $1$: for $d_{1,h}\Delta \gg 1$, $\widehat M_h \approx z_h(j)$ and no memory is left.

\subsection{The objective: the coupling in the latent space}\label{sec:objective}

The free-running fit alone does not recover the network: on the network FN8 (10 runs, on an earlier version of the code), a network learned from scratch with this fit reached an AUPR of $0.45$--$0.57$, against $0.75$ for \texttt{fit\_network}. A Sinkhorn divergence between simulated and observed cells only constrains marginals, whereas once an ancestor is chosen, the network term of the coupling requires the sigmoid to explain each gene of the cell. The objective therefore keeps the coupling, in the latent space.

For an interval $(t_a, t_b)$ of length $\Delta$, with states $i$ at $t_a$ and cells $j$ at $t_b$, the cost of a pair is
\begin{equation}\label{eq:cost}
c_{ij} = \underbrace{\sum_g w_g\, \CE\big(\sigma_g(\hat z_{ij}), \tau_{gj}\big)}_{\text{network, all genes}}
+ \underbrace{\sum_{h} \frac12 \Big[\frac{\big(z_h(j) - m_{h}\big)^2}{v_h} + \log v_h\Big]}_{\text{reachability, hidden variables}},
\end{equation}
where $\hat z_{ij}$ gathers the proteins of the anchored regulators at the end of their bridges from $i$ and the hidden states $z_h(j)$, and $m_h$, $v_h$ are given by~\eqref{eq:transition} with $\sigma_h = \sigma_h(\hat z_{ij})$. The hidden states $z_h(i)$ and $z_h(j)$ are read by an encoder on the real cell behind the state $i$ and on the cell $j$. The objective, for each interval, is
\begin{equation}\label{eq:objective}
L = \min_{\pi} \ \langle \pi, c \rangle - \varepsilon H(\pi),
\end{equation}
with $\pi$ at the marginals of the masses. It has the following properties.

\paragraph{CardamomOT is a special case.} Without hidden variables, the second term vanishes and the first is the network term $B$ of~\eqref{eq:current-cost}, on all the genes.

\paragraph{Gradient.} By the envelope theorem, the gradient of~\eqref{eq:objective} with respect to the parameters is $\langle \pi^\star, \partial c \rangle$ at the optimal coupling. There is no differentiation through the Sinkhorn iterations, and the gradient can be estimated on pairs drawn from $\pi^\star$: the loop becomes an online EM.

\paragraph{Bound.} The entropic OT with any cost bounds a mixture likelihood.

\begin{prop}\label{prop:bound}
Let $c$ be a cost between $N$ source states and $M$ target cells, and $p_{ij} = e^{-c_{ij}}$. For any coupling $\pi$ with uniform target marginal,
\[
-\frac{1}{M}\sum_j \log \Big(\frac{1}{N}\sum_i p_{ij}\Big) \ \le\ \langle \pi, c \rangle - H(\pi) + \log (MN),
\]
with equality for $\pi(i \mid j) \propto p_{ij}$.
\end{prop}

\begin{proof}
Jensen's inequality for each $j$ with $q(i) = \pi(i \mid j)$, then averaging over $j$ with $\pi_{ij} = q_j(i)/M$.
\end{proof}

With the cost~\eqref{eq:cost}, $p_{ij}$ is a product over the genes of $\sigma^{\tau}(1 - \sigma)^{1-\tau}$ raised to the weight $w_g$, times the Gaussian density of the reachability, up to factors that only depend on $j$. At $\varepsilon = 1$, the objective is therefore a bound on a pseudo-likelihood of the targets $\tau$ of the cells, the targets being themselves estimated from the counts by the mixture. It is not a likelihood of the counts. With $\varepsilon \neq 1$, it is the same bound for the tempered cost $c/\varepsilon$; the regularisation of the current trajectories is relative to the cost. The source marginal only restricts $\pi$, so the bound stays valid, and it adds a constraint of conservation of mass that the bound per cell does not contain: no state can be left unused while another explains every cell.

\paragraph{Why the cross-entropy.} The log-likelihood ratio of the NB between two modes grows with the expression of the gene, so that a few highly expressed genes would decide the coupling; with thousands of genes, the imbalance would grow. In the cross-entropy, each gene contributes a term of the same scale, whose weight is set by $w_g$. The NB likelihood has only been used in the model learning, never in the network fit, and nothing indicates that it would do better there.

\subsection{The extensions of CardamomOT as components}\label{sec:extensions}

\begin{center}
\begin{tabular}{@{}p{0.27\textwidth}p{0.25\textwidth}p{0.4\textwidth}@{}}
\toprule
Extension & Component & Point of attention \\
\midrule
Stimulus, schedules & dynamics (exogenous part of $z$) & unchanged \\
Proliferation, dilution & dynamics, masses of the coupling & marker genes are part of the data rather than a selection; in simulations, the rates must be functions of $z$ \\
Model learning of $\sigma$ & decoder, family per target & little information on weakly expressed targets \\
mRNA memory & decoder & $\sigma$ evaluated at several times for every target \\
Network conditions, genes per condition & $\theta$ per condition & gene masks may become unnecessary; the core may differ between conditions \\
Mixed effects & prior on $\theta$ & does not address the main source of differences between replicates \\
Multiome & second layer of the decoder & not one more hidden variable \\
CardamomOT-stat & same cost, self-coupling & the $d_1$ of a hidden variable is probably harder to identify without time series \\
Observed proteins & measurement of $z$ & the strongest anchor \\
Masses, transitions, lineage & constraints on the coupling & unchanged \\
\bottomrule
\end{tabular}
\end{center}

\paragraph{Multiome.} In CardamomOT-multiome, the accessibility of regulatory regions gates the regulated part of the burst frequency of their genes, and is itself regulated by the proteins through a second network. Regions are many and specific to genes: they are local states of the decoder, regulated by $z$, rather than components of $z$.

\paragraph{Mixed effects.} In CardamomOT-ME, no method recovers the deviations of the patients, and the largest source of differences between replicates is the mixture of each sample. The framework does not change this by itself; bringing the mixture into the objective (Section~\ref{sec:stage2}) is what could.

\paragraph{Perturbations.} A knock-out of a gene outside the core only changes that gene. This is a direct consequence of the framework, and it matters for Perturb-seq data, in which many perturbed genes are outside any core.

\paragraph{Proliferation as a hidden variable.} An idea to test: the proliferative state as a hidden variable, read on its signature genes and driving the dilution. It would then stop being a nuisance factor and enter the dynamics.

\paragraph{A reserve.} This unity is conceptual. The code assumes a square $\theta$ and arrays of states by genes, so that the move to all the genes touches every extension. The first stage is developed on the base model with the extensions off; they are then switched back on one at a time, checking that their results hold.

%======================================================================
\section{Stage 1: CardamomOT on all the genes}\label{sec:stage1}
%======================================================================

\paragraph{Candidates and core.} A prior gives a wide set of $C$ candidate regulators: expressed transcription factors, chromatin regulators, the regulators of the literature prior. The current gene selection only initialises the core. The core is the support of $\theta$: the candidates with at least one outgoing edge.

\paragraph{One iteration.}
\begin{enumerate}[nosep]
\item Coupling: every gene enters the network term of the cost~\eqref{eq:current-cost}; the proteins of each pair are only computed for the regulators of the core.
\item Trajectories: on the pairs drawn from the coupling, the proteins of all $C$ candidates are computed. The cost is $n\,C$ rather than $n^2 C$, since the end states are already recomputed for the drawn pairs only.
\item Fit: each target is regressed on the candidates, after a screening step~\cite{fan2008sure}, with the QUT penalty.
\item Core: the support of $\theta$ is the core of the next iteration. A candidate left out can come back, since its protein is still computed.
\end{enumerate}
The selection of the genes is then part of the inference.

\paragraph{Computing time.} Per pair, the cost requires $K$ bridges and $\sum_g \nnz(\theta_g) \le GK$ terms of the sigmoids, instead of $G$ bridges and $G^2$ terms. The main code to write is a sparse sigmoid: the input of $\sigma$ currently runs over every regulator for each target. The mixture is fitted gene by gene, linearly in $G$.

\paragraph{Weights of the satellites.} The fits are independent per target and the basins are updated gene by gene, so the satellites only act on the core through the coupling. Three consequences:
\begin{itemize}[nosep]
\item satellites enter the network term $B$ of the cost, and neither the mRNA nor the protein terms;
\item the blocks must be normalised, otherwise the satellites overwhelm the mRNAs of the core;
\item a satellite that depends on something outside the core pulls the trajectories in a wrong direction. Its weight in $B$ follows how well the core explains it (gain of likelihood against a model without regulators). Blocks of such satellites with correlated residuals show what the core is missing.
\end{itemize}

\paragraph{Lock-in.} The current core shapes the trajectories, and the trajectories can favour the genes that resemble it. This is the main risk of the first stage, and the reason for the tests of Section~\ref{sec:validation}: a true regulator removed from the initial core must come back, decoys must not stay, and the satellites must improve the core.

%======================================================================
\section{Stage 2: CardaVI}\label{sec:stage2}
%======================================================================

\paragraph{A single objective.} The aim is a single objective, of which every step of the inference is a descent, rather than a fully differentiable inference: gradients are only a means, and gradients on marginals alone do not learn the network (Section~\ref{sec:objective}). CardaVI can keep an OT solved separately, and possibly the fits per target, as long as each block descends the objective~\eqref{eq:objective}. The objective can be followed along the iterations of the first stage already: the steps that do not descend it are the inconsistencies to resolve.

\paragraph{Steps outside the objective.} The note on model learning shows that the basin update and the network fit maximise the same likelihood. The following steps are not yet part of it:
\begin{itemize}[nosep]
\item the mixture, fitted before the inference on the counts alone. Keeping its levels and rates fixed is the simplest choice; the weights of its modes per sample are a candidate for the objective (Section~\ref{sec:extensions}, mixed effects);
\item the mRNA and protein terms of the trajectory cost, and its regularisation relative to the cost;
\item the forcing of the trajectories by mass-constrained basins, rather than by the posteriors;
\item the free-running refinement, which fits a Sinkhorn divergence on marginals;
\item the reading of $\sigma$: with two modes, the basin update reads $\sigma$ as the probability of the high mode, whereas the simulator reads it as a burst frequency between the two levels. The two readings coincide when the proteins put $\sigma$ near $0$ or $1$; they differ when the proteins do not separate the cells of different modes, as along deterministic ODE trajectories.
\end{itemize}
Some of these are deliberate: the forcing by basins avoids the contraction of the expected levels. A single objective must either include them or justify them.

\paragraph{Encoder.} The encoder is a free neural network. It belongs to the inference, not to the model: its form has no mechanistic meaning, and all the structure lies in the dynamics and the decoder. Its first role is to read the hidden variables in each cell, from their targets. It also applies to new cells. It could read the anchored proteins too, which would remove the chain of ancestors and let intervals be sampled independently, at the price of the memory that the chain carries; with the chain kept, the anchored proteins are those of CardamomOT, and so is the information they carry.

\paragraph{Smooth introduction of the hidden variables.} A hidden variable starts with no outgoing edges, so that the model is exactly the model of the first stage. The weight of its edges then rises progressively, and it is kept only if it improves the score on all the data, with a penalty of BIC type. Its initial state comes from a block of satellites with correlated residuals, and a non-selected gene whose mRNA precedes the inferred state is a candidate identity, which can then be anchored.

\paragraph{Discovery of the core.} In the first stage, the core is the support of $\theta$ given by the QUT penalty. If $\theta$ is learned by gradient, gates on the outgoing edges of each candidate (a hard-concrete relaxation of an $L_0$ penalty~\cite{louizos2018learning}) play the same role. Neither lifts the cap of Section~\ref{sec:cap}: among collinear candidates the choice is arbitrary. What can be claimed is a class of regulators, stable over several initialisations, rather than one gene.

\paragraph{When to move to the second stage.} The second stage is worth its cost if the first shows at least one of the following: blocks of correlated residuals reveal missing regulators that the first stage cannot capture; the chain of ancestors blocks at scale; an inconsistency between steps degrades the results. If the satellites do not improve the core, the premise of both stages falls.

%======================================================================
\section{Identifiability and failure modes}\label{sec:risks}
%======================================================================

\paragraph{Identifiability of the latent variables.} Anchors (the mRNA of the regulator, the literature prior as a mask of the edges), sparsity and the non-linearity of $\sigma$ constrain the latent space. Several theoretical frameworks may apply: sparse decoding with anchor features~\cite{moran2022identifiable}, sparsity of the latent mechanism~\cite{lachapelle2022disentanglement}, auxiliary variables~\cite{khemakhem2020variational}, here the time and the sample, and interventions~\cite{zhang2023identifiability}. Whether a latent PDMP with a sigmoid burst frequency fits in one of them is an open result. Nuisance factors are the mechanistic analogue of the low-rank part of sparse plus low-rank models~\cite{chandrasekaran2012latent}.

\paragraph{Failure modes.}
\begin{enumerate}[nosep]
\item \emph{Lock-in of the core} (first stage, Section~\ref{sec:stage1}).
\item \emph{Start.} The encoder and the coupling depend on each other. As in the current loop, the coupling starts on the mRNA term, with the encoder frozen and the hidden variables closed, then moves to the network term.
\item \emph{Hidden variables that cheat.} The network term does not constrain the dynamics of $z_h$, since $\hat z_{ij,h} = z_h(j)$ does not depend on $i$: everything rests on the reachability term. Flexible variables may also take over the signal of the anchored regulators.
\item \emph{Vanishing noise.} In~\eqref{eq:cost}, a learned $a_h \to 0$ with a constant $z_h$ sets the first term of the reachability to zero while $\log v_h \to -\infty$: the cost is unbounded below. It takes a lower bound on $a_h$ or a stochastic encoder. Conversely, removing $\log v$ would make a large $a_h$ free.
\item \emph{Decorative dynamics.} A good reconstruction with a weak dynamics gives a variational autoencoder with a mechanistic costume. The model must only be judged on free-running simulation and on perturbations, never on reconstruction.
\end{enumerate}

%======================================================================
\section{Validation plan}\label{sec:validation}
%======================================================================

The simulated data come from the simulator of the package: a known core network, with satellites regulated by it. Each level answers one question, with one main measurement.

\begin{center}
\begin{tabular}{@{}lp{0.42\textwidth}p{0.38\textwidth}@{}}
\toprule
Level & Question & Measurement \\
\midrule
1.1 & Does using all the genes keep the network of the core? & AUPR of the core and time, against CardamomOT on the right genes only \\
1.2 & Is the core discovered? & true regulators kept and decoys removed, from a prior with decoys \\
1.3 & Does the core escape lock-in? & a true regulator removed from the initial core comes back \\
1.4 & Do the satellites improve the core? & recovery of the core against their number \\
2.1 & Do the steps descend the same objective? & the objective along the iterations of the first stage \\
2.2 & Can a hidden regulator be recovered? & its protein and its regulators, its mRNA removed \\
2.3 & Does the model predict at scale? & free-running simulation, perturbations \\
\bottomrule
\end{tabular}
\end{center}

Level 1.4 tests the central hypothesis, and level 2.2 is the decisive one for hidden variables. Level 2.3 includes Perturb-seq data, on which deep models do not outperform simple linear baselines~\cite{ahlmann2025deep}: a mechanistic model at the scale of the transcriptome would have a clear benchmark. A prospective validation, with perturbations predicted before they are tested, would complete it.

\tofill{results.}

%======================================================================
\section{Open questions}\label{sec:open}
%======================================================================

\begin{itemize}[nosep]
\item Which encoder: a single cell, or a recursive form $E(x_j, m(i,j))$ that also receives the state predicted from the ancestor, closer to a filter?
\item Learning the network by gradient: the QUT penalty is calibrated for fits per target. The first versions keep these fits.
\item The identifiability result of Section~\ref{sec:risks}.
\item The Gaussian transition at frozen $\sigma_h$ over long intervals: how many integration steps are needed?
\item Does the difference between the two readings of $\sigma$ matter in practice?
\item Does the dispersion signature of the direction of a switch survive in real data?
\item Minibatch couplings are biased: how much, at the scale of the transcriptome?
\end{itemize}

\begin{thebibliography}{99}

\bibitem{mauge2026cardamomot}
Y.~Maugé, E.~Ventre.
\newblock CardamomOT: a mechanistic optimal transport-based framework for gene regulatory network inference, trajectory reconstruction and generative modeling.
\newblock \emph{bioRxiv}, 2026. doi:10.64898/2026.03.31.715390. Accepted in \emph{PLoS Computational Biology}.

\bibitem{herbachinferring2017}
U.~Herbach, A.~Bonnaffoux, T.~Espinasse, O.~Gandrillon.
\newblock Inferring gene regulatory networks from single-cell data: a mechanistic approach.
\newblock \emph{BMC Systems Biology}, 11:105, 2017.

\bibitem{lopez2018deep}
R.~Lopez, J.~Regier, M.~B.~Cole, M.~I.~Jordan, N.~Yosef.
\newblock Deep generative modeling for single-cell transcriptomics.
\newblock \emph{Nature Methods}, 15:1053--1058, 2018.

\bibitem{carilli2024biophysical}
M.~Carilli, G.~Gorin, Y.~Choi, T.~Chari, L.~Pachter.
\newblock Biophysical modeling with variational autoencoders for bimodal, single-cell RNA sequencing data.
\newblock \emph{Nature Methods}, 2024.

\bibitem{bertin2025scalable}
P.~Bertin, J.~D.~Viviano, A.~Tejada-Lapuerta, W.~Wang, S.~Bauer, F.~J.~Theis, Y.~Bengio.
\newblock A scalable gene network model of regulatory dynamics in single cells.
\newblock \emph{arXiv:2503.20027}, 2025.

\bibitem{moss2020gene}
J.~Moss, P.~Liò.
\newblock Gene regulatory network inference with latent force models.
\newblock \emph{arXiv:2010.02555}, 2020.

\bibitem{weinreb2018fundamental}
C.~Weinreb, S.~Wolock, B.~K.~Tusi, M.~Socolovsky, A.~M.~Klein.
\newblock Fundamental limits on dynamic inference from single-cell snapshots.
\newblock \emph{Proceedings of the National Academy of Sciences}, 115(10):E2467--E2476, 2018.

\bibitem{lavenant2021towards}
H.~Lavenant, S.~Zhang, Y.-H.~Kim, G.~Schiebinger.
\newblock Towards a mathematical theory of trajectory inference.
\newblock \emph{arXiv:2102.09204}, 2021.

\bibitem{fan2008sure}
J.~Fan, J.~Lv.
\newblock Sure independence screening for ultrahigh dimensional feature space.
\newblock \emph{Journal of the Royal Statistical Society: Series B}, 70(5):849--911, 2008.

\bibitem{louizos2018learning}
C.~Louizos, M.~Welling, D.~P.~Kingma.
\newblock Learning sparse neural networks through $L_0$ regularization.
\newblock \emph{International Conference on Learning Representations}, 2018.

\bibitem{moran2022identifiable}
G.~E.~Moran, D.~Sridhar, Y.~Wang, D.~M.~Blei.
\newblock Identifiable deep generative models via sparse decoding.
\newblock \emph{Transactions on Machine Learning Research}, 2022.

\bibitem{lachapelle2022disentanglement}
S.~Lachapelle, P.~Rodríguez, Y.~Sharma, K.~Everett, R.~Le~Priol, A.~Lacoste, S.~Lacoste-Julien.
\newblock Disentanglement via mechanism sparsity regularization: a new principle for nonlinear ICA.
\newblock \emph{Conference on Causal Learning and Reasoning}, 2022.

\bibitem{khemakhem2020variational}
I.~Khemakhem, D.~Kingma, R.~Monti, A.~Hyvärinen.
\newblock Variational autoencoders and nonlinear ICA: a unifying framework.
\newblock \emph{International Conference on Artificial Intelligence and Statistics}, 2020.

\bibitem{zhang2023identifiability}
J.~Zhang, K.~Greenewald, C.~Squires, A.~Srivastava, K.~Shanmugam, C.~Uhler.
\newblock Identifiability guarantees for causal disentanglement from soft interventions.
\newblock \emph{Advances in Neural Information Processing Systems}, 2023.

\bibitem{chandrasekaran2012latent}
V.~Chandrasekaran, P.~A.~Parrilo, A.~S.~Willsky.
\newblock Latent variable graphical model selection via convex optimization.
\newblock \emph{The Annals of Statistics}, 40(4):1935--1967, 2012.

\bibitem{ahlmann2025deep}
C.~Ahlmann-Eltze, W.~Huber, S.~Anders.
\newblock Deep-learning-based gene perturbation effect prediction does not yet outperform simple linear baselines.
\newblock \emph{Nature Methods}, 2025.

\end{thebibliography}

\end{document}
